
\paragraph{NK landscapes and collective search.}
The NK landscape model~\citep{kauffman1993origins} provides a tunable framework
for studying optimization on rugged fitness landscapes, parameterized by the
number of components~$N$ and the degree of epistasis~$K$. It has been widely
applied to organizational decision-making~\citep{levinthal1997adaptation,
rivkin2000imitation}, where groups search for high-fitness configurations under
varying degrees of interdependence. In these models, organizations navigate
rugged landscapes through hierarchical decomposition, imitation, or parallel
search. Our work extends this tradition by modeling the search process as
democratic voting rather than managerial or evolutionary search, and by
introducing a partition between shared (law-like) and individual (trait-like)
components that creates voter heterogeneity absent from prior NK collective
search models.

\paragraph{Computational social choice.}
The formal study of voting methods has a long
history~\citep{arrow1951social, gibbard1973manipulation, satterthwaite1975strategy}. Arrow's impossibility
theorem established that no rank-order voting system can satisfy a small set of
reasonable criteria simultaneously, while the Gibbard--Satterthwaite theorem
showed that deterministic, non-dictatorial methods are susceptible to strategic manipulation. Computational approaches have
since analyzed the complexity, manipulability, and axiomatic properties of
various methods~\citep{brandt2016handbook}. However, most analyses focus on
single-shot elections or equilibrium outcomes rather than iterated collective
optimization over many rounds, which is the regime most relevant to ongoing
governance. Our framework bridges this gap by evaluating voting methods as
long-run optimizers on a landscape where the quality of collective decisions
accumulates over time.

\paragraph{Law as a complex adaptive system.}
Legal scholars have increasingly applied complexity theory to
law~\citep{ruhl1996fitness, ruhl2008complexity}, arguing that legal systems
exhibit emergent behavior, path dependence, and co-evolution with the populations
they govern. Our model operationalizes this perspective: laws (shared bits) and
individual traits co-determine fitness through epistatic interactions, and the
voting process drives adaptation of the legal component. The cross-dependency
parameter $\alpha$ formalizes the degree to which legislation interacts with
individual circumstances, a central concern in regulatory design that
complexity-theoretic analyses of law have identified but not previously modeled
at the level of voting mechanisms.
